% =========================================================================
\section{Circuito 2: amplificador no inversor}\label{sec:no-inversor}
% =========================================================================
Montamos el no inversor de la \cref{fig:no-inversor}, con $1+R_2/R_1=11$, a $1\,\mathrm{kHz}$.

% Figura: amplificador no inversor (circuito 2) con tensiones y corrientes del análisis.
\begin{figure}[ht]
\centering
\begin{circuitikz}[circuito]
  \draw (0,0) node[op amp, ao, noinv input up] (A) {};
  % nudo inversor (abajo) y red de realimentación
  \draw (A.out) -- ++(0.7,0) coordinate (O);
  \draw (A.-) -- ++(-0.6,0) coordinate (N) -- ++(0,-1.0) coordinate (M);
  \draw (M) -- ++(-1.0,0) coordinate (G1) to[R, l_=$R_1$, i>^=$i_1$] ++(0,-1.7) node[ground] {};
  \draw (M) to[R, l_=$R_2$, i<_=$i_2$] (M -| O) -- (O);
  % generador en la entrada no inversora
  \draw (A.+) to[short, -o] ++(-2.0,0) coordinate (Ein) node[above, tens] {$\vin$};
  \draw (Ein) -- ++(-0.8,0) to[sV, l_=$v_g$] ++(0,-2.0) node[ground] {};
  % carga y salida
  \draw (O) to[short, *-o] ++(1.0,0) node[right, tens] {$\vout$};
  \draw (O) ++(0.5,0) to[R, l=$R_L$] ++(0,-2.0) node[ground] {};
  \fill (M) circle (1.6pt);
  \node[tens, anchor=east, font=\scriptsize] at ($(N)!0.5!(M)+(-0.08,0)$) {$\vm=\vin$};
\end{circuitikz}
\caption{Amplificador no inversor (figura~3 del guion), con $R_1=1\,\mathrm{k}\Omega$,
$R_L=10\,\mathrm{k}\Omega$ y la alimentación omitida. La entrada $+$ no absorbe corriente; el nudo
inversor sigue a la entrada ($\vm=\vin$) y $R_1$--$R_2$ forman un divisor de $\vout$.}
\label{fig:no-inversor}
\end{figure}


\begin{pregunta}[Formas de onda y ganancia]
Dibujar las formas de onda de $\Vin=V_g$ y $\Vout$ detallando las escalas, y comparar la relación de
amplitudes o tensiones pico a pico con la $A_V$ buscada.
\end{pregunta}

\begin{respuesta}
Ahora la ganancia teórica y la experimental son
\begin{equation}\label{eq:n-Av}
  A_V^{\mathrm{teo}}=1+\frac{R_2}{R_1},\qquad
  A_V^{\mathrm{exp}}=+\frac{V_{\mathrm{out},pp}}{V_{\mathrm{in},pp}} ,
\end{equation}
con signo positivo porque la salida está en fase con la entrada.
\begin{center}\small
\renewcommand{\arraystretch}{1.3}
\begin{tabular}{@{}cccccc@{}}
\toprule
$V_p$ & Escala X, Y (V/div) & Base de tiempos & $V_{\mathrm{in},pp}$ & $V_{\mathrm{out},pp}$ & $A_V^{\mathrm{exp}}$\\
\midrule
\dato & \dato & \dato & \dato & \dato & \dato\\
\dato & \dato & \dato & \dato & \dato & \dato\\
\bottomrule
\end{tabular}
\end{center}
\begin{center}
\hueco{0.6\textwidth}{Oscilograma $\vin$, $\vout$ (en fase)}
\end{center}
Obtenemos $A_V^{\mathrm{exp}}=$ \dato\ frente a $A_V^{\mathrm{teo}}=$ \dato, una desviación de \dato. A diferencia del inversor,
la entrada no absorbe corriente ($\Zin\to\infty$), de modo que $\vin$ coincide con la tensión del generador.
\end{respuesta}

\begin{pregunta}[Característica de transferencia y recorte]
Dibujar la característica $\Vout(V_g)$. Señalar para qué amplitud de entrada (en la representación temporal) se
produce el recorte y relacionarlo con lo medido en la representación X--Y.
\end{pregunta}

\begin{respuesta}
La característica es una recta que pasa por el origen con pendiente $+11$ (la ganancia), hasta que la salida
llega a la tensión de saturación $\Vsat$; a partir de ahí la salida ya no puede crecer y la curva se vuelve
horizontal.

La salida empieza a recortarse cuando $11\cdot V_p$ alcanza $\Vsat$, es decir, para una amplitud de entrada
\begin{equation}\label{eq:n-umbral}
  V_{p,\mathrm{umbral}}=\frac{\Vsat}{A_V}=\frac{\Vsat}{11} .
\end{equation}
Medimos $V_{p,\mathrm{umbral}}=$ \dato\ en la pantalla en función del tiempo, y en el modo X--Y los codos de la
característica aparecen en $\vin=$ \dato\ y \dato. Ambas medidas deben coincidir, porque son el mismo
fenómeno visto de dos formas: el punto en que la recta se dobla en el modo X--Y es justo la amplitud a la
que la sinusoide de salida empieza a aplanarse.
\begin{center}
\hueco{0.6\textwidth}{Característica $\vout(\vin)$ en modo X--Y}
\end{center}
\end{respuesta}

\begin{pregunta}[Definición de la ganancia]
A la vista de las características de transferencia de ambos circuitos, discutir si sería conveniente o
posible definir la ganancia como $A_V=\partial\Vout/\partial\Vin$.
\end{pregunta}

\begin{respuesta}
Las dos características son lineales a trozos: una recta por el origen en la zona lineal y tramos
horizontales en saturación. La definición como cociente, $A_V=\vout/\vin$, solo tiene sentido en la zona
lineal; en saturación el cociente varía con la amplitud de entrada y no describe el circuito. La
definición incremental,
\begin{equation}\label{eq:n-deriv}
  A_V=\frac{\partial\vout}{\partial\vin}=
  \begin{cases}
    -R_2/R_1 \ \text{(inversor)},\quad 1+R_2/R_1 \ \text{(no inversor)}, & \text{zona lineal},\\[2pt]
    0, & \text{saturación},
  \end{cases}
\end{equation}
es la pendiente de la característica en cada punto. En la zona lineal coincide con el cociente, porque la
recta pasa por el origen; pero es más general: describe la respuesta a pequeñas variaciones en torno a cualquier
punto de trabajo, sigue siendo válida si hubiera una tensión de \emph{offset} (la recta no pasaría por el origen
y el cociente variaría con $\vin$) y muestra explícitamente que en saturación el amplificador deja de
amplificar. Por eso sí es conveniente, y es la definición que se usa para el comportamiento en pequeña señal.
\end{respuesta}
